\section*{अध्याय \ref{ch:b1:matrices} --- आव्यूह}

\begin{solution}{exo:b1:matrices:1}
\[
AB = \begin{pmatrix} 2 & 1\\ 1 & 0\end{pmatrix},
\quad
BA = \begin{pmatrix} 0 & 1\\ 1 & 2\end{pmatrix},
\quad
A^2 - B^2 = \begin{pmatrix} 1 & 4\\ 0 & 1\end{pmatrix} - I
= \begin{pmatrix} 0 & 4\\ 0 & 0\end{pmatrix},
\]
\[
(A+B)(A-B) = A^2 - AB + BA - B^2
= \begin{pmatrix} 0 & 4\\ 0 & 0\end{pmatrix} +
\begin{pmatrix} -2 & 0\\ 0 & 2 \end{pmatrix}
= \begin{pmatrix} -2 & 4\\ 0 & 2\end{pmatrix}.
\]
वे $BA - AB \neq 0$ जितने भिन्न हैं: सर्वसमिका $(a+b)(a-b) = a^2 - b^2$ के
लिए क्रमविनिमयता चाहिए, जो यहाँ विफल है।
\end{solution}

\begin{solution}{exo:b1:matrices:2}
$M$: $L_2 \leftarrow L_2 - 2L_1$ दूसरी पंक्ति मार देता है; $L_3 \leftarrow
L_3 - L_1$ $(0, -1, -2)$ देता है। दो धुरियाँ: $\operatorname{rk} M = 2$।

$N$: $L_3 \leftarrow L_3 - L_1$ $(0,1,1,1) = L_2$ देता है; फिर $L_3
\leftarrow L_3 - L_2 = 0$। दो धुरियाँ: $\operatorname{rk} N = 2$।
\end{solution}

\begin{solution}{exo:b1:matrices:3}
$(A \mid I_3)$ का लघुकरण: $L_2 \leftarrow L_2 - 2L_1$, $L_3 \leftarrow L_3 -
L_1$:
\[
\begin{pmatrix}
1 & 0 & 1 & 1 & 0 & 0\\
0 & 1 & -1 & -2 & 1 & 0\\
0 & 1 & 0 & -1 & 0 & 1
\end{pmatrix}
\xrightarrow{L_3 \leftarrow L_3 - L_2}
\begin{pmatrix}
1 & 0 & 1 & 1 & 0 & 0\\
0 & 1 & -1 & -2 & 1 & 0\\
0 & 0 & 1 & 1 & -1 & 1
\end{pmatrix},
\]
फिर $L_1 \leftarrow L_1 - L_3$, $L_2 \leftarrow L_2 + L_3$:
\[
A^{-1} = \begin{pmatrix}
0 & 1 & -1\\
-1 & 0 & 1\\
1 & -1 & 1
\end{pmatrix}.
\]
\emph{जाँच:} $A$ की पहली पंक्ति गुणा $A^{-1}$ का पहला स्तंभ: $1 \cdot 0 +
0\cdot(-1) + 1\cdot 1 = 1$; गुणा दूसरा स्तंभ: $1 - 0 - 1 = 0$; गुणा तीसरा:
$-1 + 0 + 1 = 0$।
\end{solution}

\begin{solution}{exo:b1:matrices:4}
$u(1) = 1$, $u(X) = X + 1$, $u(X^2) = X^2 + 2X + 1$: $(1, X, X^2)$ में
निर्देशांकों के स्तंभ देते हैं
\[
M = \begin{pmatrix}
1 & 1 & 1\\
0 & 1 & 2\\
0 & 0 & 1
\end{pmatrix}.
\]
$u$ व्युत्क्रमणीय है, क्योंकि उसका स्पष्ट प्रतिलोम $P \mapsto P(X - 1)$ है
(प्रतिस्थापनों का संयोजन)। उसका आव्यूह उसी तरह $u^{-1}(X^k) = (X-1)^k$ से
मिल जाता है:
\[
M^{-1} = \begin{pmatrix}
1 & -1 & 1\\
0 & 1 & -2\\
0 & 0 & 1
\end{pmatrix}.
\]
\end{solution}

\begin{solution}{exo:b1:matrices:5}
$N = \begin{pmatrix} 0 & 1\\ 0 & 0\end{pmatrix}$, $N^2 = 0$। चूँकि $2I$ और
$N$ क्रमविनिमेय हैं, द्विपद प्रसार दो पदों के बाद कट जाता है:
\[
A^k = (2I + N)^k = 2^k I + k\,2^{k-1} N
= \begin{pmatrix} 2^k & k\,2^{k-1}\\ 0 & 2^k\end{pmatrix}.
\]
($k = 2$ पर जाँच: $A^2 = \begin{pmatrix}4 & 4\\ 0 & 4\end{pmatrix}$, जो सीधे
गुणनफल से सही है।)
\end{solution}

\begin{solution}{exo:b1:matrices:6}
\omterm{def:b1:matrices:transpose}{अनुरेख}: $\operatorname{tr}(AB - BA) = \operatorname{tr}(AB) -
\operatorname{tr}(BA) = 0$ (\cref{def:b1:matrices:transpose}), जबकि $\R$
अथवा $\C$ में $\operatorname{tr}(I_n) = n \neq 0$। अतः कोई हल नहीं।
(अनंत-विमीय समष्टियों पर यह सर्वसमिका साकार \emph{होती है} --- अवकलन और $x$
से गुणन उसे संतुष्ट करते हैं --- और ठीक इसीलिए, क्योंकि वहाँ कोई \omterm{def:b1:matrices:transpose}{अनुरेख} होता
ही नहीं।)
\end{solution}

\begin{solution}{exo:b1:matrices:7}
दूरबीनी गुणनफल, क्योंकि $A$ की सभी घातें क्रमविनिमेय हैं:
\[
(I - A)(I + A + \dots + A^{m-1}) = I - A^m = I ,
\]
और \cref{prop:b1:matrices:ring} एकतरफ़ा प्रतिलोम को उन्नत कर देता है।
अनुप्रयोग के लिए: दिया हुआ आव्यूह $I + N$ है, जहाँ
\[
N = \begin{pmatrix} 0 & 2 & 3\\ 0 & 0 & 2\\ 0&0&0 \end{pmatrix},
\quad
N^2 = \begin{pmatrix} 0&0&4\\ 0&0&0\\ 0&0&0\end{pmatrix},
\quad N^3 = 0 ,
\]
अतः सूत्र में $A$ के स्थान पर $-N$ रखने पर:
\[
(I + N)^{-1} = I - N + N^2 =
\begin{pmatrix}
1 & -2 & 1\\
0 & 1 & -2\\
0 & 0 & 1
\end{pmatrix}.
\]
\end{solution}

\begin{solution}{exo:b1:matrices:8}
$A^2 = A$: अंतःसमाकारिता $a$ एक \omterm{def:b1:linmaps:projection}{प्रक्षेप} है
(\cref{thm:b1:linmaps:projchar}), $\dim\operatorname{im} a = r =
\operatorname{rk} A$ के साथ $E = \operatorname{im} a \oplus \ker a$। इस
विघटन के अनुकूल किसी \omterm{def:b1:vspaces:free}{आधार} में (\omterm{def:b1:linmaps:kerim}{प्रतिबिंब} के $r$ सदिश, फिर अष्टि का \omterm{def:b1:vspaces:free}{आधार}) $a$
का आव्यूह $\begin{pmatrix} I_r & 0\\ 0 & 0\end{pmatrix}$ है, जिसका \omterm{def:b1:matrices:transpose}{अनुरेख}
$r$ है। \omterm{def:b1:matrices:transpose}{अनुरेख} \omterm{def:b1:matrices:changeofbasis}{आधार-परिवर्तन} के अंतर्गत अपरिवर्तित रहता है: चक्रीय सर्वसमिका
से $\operatorname{tr}(P^{-1}MP) = \operatorname{tr}(MPP^{-1}) =
\operatorname{tr} M$। अतः $\operatorname{tr} A = r = \operatorname{rk} A$।
\end{solution}

\begin{solution}{exo:b1:matrices:9}
$J^2 = nJ$ ($J^2$ की हर प्रविष्टि $n$ इकाइयों का योग है)। $M^{-1} = \alpha I
+ \beta J$ खोजिए:
\[
(aI + bJ)(\alpha I + \beta J)
= a\alpha\, I + (a\beta + b\alpha + nb\beta)\, J .
\]
यह $I$ के बराबर है तभी जब $a\alpha = 1$ और $a\beta + b\alpha + nb\beta = 0$,
अर्थात् $\alpha = \frac1a$ तथा $\beta(a + nb) = -\frac ba$। यदि $a \neq 0$
और $a + nb \neq 0$:
\[
M^{-1} = \frac 1a I - \frac{b}{a(a + nb)}\, J .
\]
विलोमतः, यदि $a = 0$: तो $M = bJ$ की कोटि $\leq 1 < n$ है ($n \geq 2$ के
लिए): अतः व्युत्क्रमणीय नहीं ($n = 1$ अदिश स्थिति है)। यदि $a + nb = 0$: तो
सदिश $v = (1, \dots, 1)^{\mathsf T}$ $v \neq 0$ के साथ $Mv = (a + nb)v = 0$
संतुष्ट करता है: अतः व्युत्क्रमणीय नहीं। इस प्रकार $M \in GL_n \iff a \neq
0$ और $a + nb \neq 0$।
\end{solution}

\begin{solution}{exo:b1:matrices:10}
\emph{योग:} $\operatorname{im}(A + B) \subseteq \operatorname{im} A +
\operatorname{im} B$ (हर $(A+B)x = Ax + Bx$), और ग्रासमान किसी योग की विमा
को विमाओं के योग से परिबद्ध कर देता है।

\emph{सिल्वेस्टर:} मान लीजिए $a$ $V = \operatorname{im} B$ तक सीमित $A$ का
\omterm{def:b1:logic:map}{प्रतिचित्रण} है (विमा $\operatorname{rk} B$)। उसका \omterm{def:b1:linmaps:kerim}{प्रतिबिंब}
$\operatorname{im}(AB)$ है ($a(Bx) = ABx$), और $V$ में कोटि--शून्यता:
\[
\operatorname{rk} B = \dim\ker(a_{|V}) + \operatorname{rk}(AB) .
\]
अब $\ker(a_{|V}) \subseteq \ker A$, जिसकी विमा $n - \operatorname{rk} A$ है:
अतः
\[
\operatorname{rk}(AB) \geq \operatorname{rk} B - (n -
\operatorname{rk} A) = \operatorname{rk} A + \operatorname{rk} B -
n . \qedhere
\]
\end{solution}

\begin{solution}{exo:b1:matrices:11}
\begin{enumerate}
  \item प्रविष्टि-दर-प्रविष्टि, $(AD)_{ij} = a_{ij}\,d_j$ और $(DA)_{ij} =
  d_i\,a_{ij}$। अतः $AD = DA$ तभी जब सभी $i, j$ के लिए $a_{ij}(d_j - d_i) =
  0$; और $i \neq j$ होने पर गुणनखंड $d_j - d_i$ अशून्य है, जो $a_{ij} = 0$
  बाध्य कर देता है: अर्थात् $A$ विकर्ण है। विलोमतः विकर्ण आव्यूह आपस में
  क्रमविनिमेय हैं।
  \item यदि $A$ हर आव्यूह के साथ क्रमविनिमेय है, तो वह
  $\operatorname{diag}(1, 2, \dots, n)$ के साथ भी है, अतः (1) से $A =
  \operatorname{diag}(\lambda_1, \dots, \lambda_n)$। तब $A E_{ij} =
  \lambda_i E_{ij}$ ($E_{ij}$ की केवल $i$ पंक्ति बचती है) जबकि $E_{ij} A =
  \lambda_j E_{ij}$: अतः $E_{ij}$ के साथ क्रमविनिमय $\lambda_i = \lambda_j$
  बाध्य कर देता है। इस प्रकार $A = \lambda I_n$; और अदिश आव्यूह सचमुच हर
  चीज़ के साथ क्रमविनिमेय हैं। $\mathcal{M}_n(K)$ का केंद्र $K\,I_n$ है।
\end{enumerate}
\end{solution}

\begin{solution}{exo:b1:matrices:12}
\begin{enumerate}
  \item यदि $\operatorname{rk} A = 1$: तो $A$ का \omterm{def:b1:linmaps:kerim}{प्रतिबिंब} एक रेखा
  $\operatorname{Vect}(C)$ है, $C \neq 0$, अतः $A$ का $j$-वाँ स्तंभ अदिशों
  $\ell_j$ (सब शून्य नहीं) के लिए $\ell_j\,C$ है, अर्थात् $L = (\ell_1,
  \dots, \ell_n) \neq 0$ के साथ $A = C L$। विलोमतः यदि $A = CL \neq 0$, तो
  सभी स्तंभ $C$ के गुणज हैं: कोटि $1$।
  \item $A^2 = C\,(L C)\,L$, और $LC$ अदिश $\sum_i \ell_i c_i =
  \operatorname{tr}(CL) = \operatorname{tr} A$ है। अतः $A^2 =
  (\operatorname{tr} A)\,A$, इसलिए आगमन से $A^m = (\operatorname{tr}
  A)^{m-1} A$। यदि $\operatorname{tr} A \neq 0$, तो कोई भी घात शून्य नहीं
  होती; और यदि $\operatorname{tr} A = 0$, तो $A^2 = 0$: अर्थात् एक-कोटि
  आव्यूह शून्यंभावी है तभी जब उसका \omterm{def:b1:matrices:transpose}{अनुरेख} शून्य हो।
  \item $t = \operatorname{tr} A \neq -1$ के साथ:
        \[
        (I_n + A)\Bigl(I_n - \frac{A}{1 + t}\Bigr)
        = I_n + A - \frac{A + A^2}{1 + t}
        = I_n + A - \frac{(1 + t)A}{1 + t} = I_n ,
        \]
        $A^2 = tA$ का प्रयोग करते हुए। यदि $t = -1$: तो $A \neq 0$ के साथ
        $(I_n + A)A = A + A^2 = A - A = 0$, अतः $I_n + A$ $A$ के हर (अशून्य)
        स्तंभ को मार देता है: अतः न \omterm{def:b1:logic:inj}{एकैकी}, न व्युत्क्रमणीय।
\end{enumerate}
\end{solution}

\begin{solution}{pb:b1:matrices:1}
\textbf{1.} $X^n$ का घात-$2$ वाले \omterm{def:b1:poly:def}{इकाई-अग्र} $D$ से यूक्लिडीय भाग
(\cref{thm:b1:poly:division}): भागफल और शेषफल हैं और अद्वितीय हैं, तथा शेषफल
की घात $\leq 1$ है: $X^n = Q_n D + a_n X + b_n$। $n = 0$ के लिए: $Q_0 = 0$,
$(a_0, b_0) = (0, 1)$; $n = 1$ के लिए: $(a_1, b_1) = (1, 0)$।

\textbf{2.} $X$ से गुणा कीजिए और $X^2 = D + sX - p$ का लघुकरण कीजिए:
\[
X^{n+1} = X Q_n D + a_n X^2 + b_n X
= (X Q_n + a_n)\,D + (s\,a_n + b_n)\,X - p\,a_n .
\]
अंतिम व्यंजक शेषफल के आकार का है (घात $\leq 1$), अतः अद्वितीयता से $a_{n+1}
= s a_n + b_n$ और $b_{n+1} = -p a_n$। $a_{n+2} = s a_{n+1} + b_{n+1}$ में
$b_{n+1} = -pa_n$ प्रतिस्थापित करने पर $a_{n+2} = s\,a_{n+1} - p\,a_n$ मिलता
है।

\textbf{3.} $X^n = Q_n D + a_n X + b_n$ का मूलों पर मूल्यांकन कीजिए:
$\lambda^n = a_n\lambda + b_n$ और $\mu^n = a_n\mu + b_n$। घटाकर $\lambda -
\mu \neq 0$ से भाग देने पर:
\[
a_n = \frac{\lambda^n - \mu^n}{\lambda - \mu},
\qquad
b_n = \lambda^n - a_n\lambda
= \frac{\lambda\mu^n - \mu\lambda^n}{\lambda - \mu} .
\]

\textbf{4.} दुहरे मूल पर: $\lambda^n = a_n\lambda + b_n$। सर्वसमिका का अवकलन
करने पर $nX^{n-1} = Q_n'\,(X - \lambda)^2 + 2Q_n\,(X - \lambda) + a_n$, और
$\lambda$ पर मूल्यांकन करने पर: $a_n = n\lambda^{n-1}$; फिर $b_n = \lambda^n
- n\lambda^{n} = (1 - n)\lambda^{n}$।

\textbf{5.} $P = \sum_i p_i X^i$ और $Q = \sum_j q_j X^j$ के लिए,
\[
P(M)\,Q(M) = \sum_{i,j} p_i q_j M^{i+j} = (PQ)(M),
\]
क्योंकि अकेले आव्यूह $M$ की घातें आपस में क्रमविनिमेय हैं (योग तो रैखिकता से
स्पष्ट हैं)। यदि $D(M) = 0$, तो $X^n = Q_n D + a_n X + b_n$ में $M$
प्रतिस्थापित करने पर $M^n = Q_n(M)\,D(M) + a_n M + b_n I = a_n M + b_n I$
मिलता है।

\textbf{6.} सीधे गुणनफल:
\[
A^2 = \begin{pmatrix}
a^2 + bc & b(a + d)\\
c(a + d) & d^2 + bc
\end{pmatrix},
\qquad
s A = \begin{pmatrix}
a(a+d) & b(a+d)\\
c(a+d) & d(a+d)
\end{pmatrix},
\]
अतः $A^2 - sA$ की विकर्णेतर प्रविष्टियाँ शून्य हैं और विकर्ण प्रविष्टियाँ
$a^2 + bc - a^2 - ad = bc - ad = -p$: $A^2 - sA + pI_2 = 0$।

\textbf{7.} $A' = \begin{pmatrix} a' & b'\\ c' & d'\end{pmatrix}$ के साथ
$p(AA') = (aa' + bc')(cb' + dd') - (ab' + bd')(ca' + dc')$ का प्रसार कीजिए:
पद $aa'cb'$ और $ab'ca'$ कट जाते हैं, पद $bc'dd'$ और $bd'dc'$ कट जाते हैं, और
जो बचता है वह है
\[
aa'dd' - bca'd' + bcb'c' - adb'c'
= (ad - bc)(a'd' - b'c') = p(A)\,p(A').
\]
यदि $p \neq 0$, तो केली--हैमिल्टन $A\,\bigl(\tfrac1p(sI_2 - A)\bigr) =
\tfrac1p(sA - A^2) = I_2$ देता है, जिससे प्रतिलोम मिल जाता है (और
\cref{prop:b1:matrices:ring} उसे दुतरफ़ा बना देता है)। यदि $p = 0$ और $A$
व्युत्क्रमणीय होता, तो गुणात्मकता $1 = p(I_2) = p(A)\,p(A^{-1}) = 0$ देती:
असंभव। अतः $A \in GL_2 \iff p \neq 0$।

\textbf{8.} $s = 3$, $p = 2$, $D = X^2 - 3X + 2 = (X - 1)(X - 2)$: $\lambda
= 2$, $\mu = 1$, अतः $a_n = 2^n - 1$ और $b_n = 2 - 2^n$ (प्रश्न 3)। इस
प्रकार
\[
A^n = (2^n - 1)A + (2 - 2^n)I
= \begin{pmatrix} 1 & 2^n - 1\\ 0 & 2^n \end{pmatrix}.
\]
जाँच: $A^2 = \begin{pmatrix} 1 & 3\\ 0 & 4\end{pmatrix}$, सूत्र से भी और
सीधे वर्ग करने से भी।

\textbf{9.} $s = 4$, $p = 3\cdot1 - 1\cdot(-1) = 4$: $D = X^2 - 4X + 4 = (X
- 2)^2$, दुहरा मूल $\lambda = 2$। प्रश्न 4: $a_n = n\,2^{n-1}$, $b_n = (1 -
n)2^n$, अतः
\[
A^n = n\,2^{n-1}A + (1 - n)2^n I
= 2^{n-1}\begin{pmatrix} n + 2 & n\\ -n & 2 - n
\end{pmatrix}.
\]
$n = 2$ पर: $2\begin{pmatrix} 4 & 2\\ -2 & 0\end{pmatrix} = \begin{pmatrix}
8 & 4\\ -4 & 0 \end{pmatrix}$, जो सीधे संगणित $A^2$ ही है।

\textbf{10.} आगमन: $F^1 = \begin{pmatrix} F_2 & F_1\\ F_1 &
F_0\end{pmatrix}$, और
\[
F^{n+1} = F^n F =
\begin{pmatrix} F_{n+1} + F_n & F_{n+1}\\
F_n + F_{n-1} & F_n \end{pmatrix}
= \begin{pmatrix} F_{n+2} & F_{n+1}\\ F_{n+1} & F_n
\end{pmatrix}.
\]
यहाँ $s = 1$, $p = -1$, $D = X^2 - X - 1$, जिसके मूल $\varphi, \psi$ हैं
($\varphi - \psi = \sqrt5$, $\varphi\psi = -1$)। अनुक्रम $(F_n)$ में $F_0 =
0 = a_0$, $F_1 = 1 = a_1$ है और वह $(a_n)$ जैसी ही पुनरावृत्ति का पालन करता
है: $F_n = a_n = (\varphi^n - \psi^n)/\sqrt5$, अर्थात् बिने का सूत्र।
कासिनी: प्रश्न 7 की गुणात्मकता को $F^n$ पर लगाने से,
\[
F_{n+1}F_{n-1} - F_n^2 = p(F^n) = p(F)^n = (-1)^n .
\]

\textbf{11.} शर्त \omterm{def:b1:linmaps:def}{रैखिक} है और उसमें शून्य अनुक्रम है: अतः एक \omterm{def:b1:vspaces:subspace}{उपसमष्टि}। आगमन
से $u_0, u_1$ $u$ को \omterm{def:b1:linmaps:forms}{रैखिक रूप} से तय कर देते हैं, और आरंभिक मानों की हर
जोड़ी ठीक एक ही हल से साकार होती है: \cref{exo:b1:findim:10} की तरह $E_D$
$(u_0, u_1) \in K^2$ के द्वारा \omterm{def:b1:logic:inj}{एकैकी आच्छादक} तथा \omterm{def:b1:linmaps:forms}{रैखिक रूप} से प्राचलित है:
$\dim E_D = 2$।

\textbf{12.} $(a_n)$ पुनरावृत्ति का पालन करता है (प्रश्न 2), जहाँ $a_0 = 0$,
$a_1 = 1$। $(b_n)$ भी वैसा ही करता है: $b_{n+2} = -p\,a_{n+1} = -p(s a_n +
b_n) = s\,b_{n+1} - p\,b_n$ ($b_{n+1} = -pa_n$ का दो बार प्रयोग करके), जहाँ
$b_0 = 1$, $b_1 = 0$। तब संयोजन $v_n = u_1 a_n + u_0 b_n$ ऐसा हल है जिसमें
$v_0 = u_0$, $v_1 = u_1$; और एक ही आरंभिक मानों वाले दो हल एक ही होते हैं
(आगमन), अतः सभी $n$ के लिए $u_n = u_1 a_n + u_0 b_n$।

\textbf{13.} $(\lambda^n)$ एक हल है तभी जब सभी $n$ के लिए $\lambda^{n+2} =
s\lambda^{n+1} - p\lambda^n$, अर्थात् $D(\lambda) = 0$ ($\lambda^n \neq 0$
से भाग देने के बाद; ध्यान दीजिए कि $\lambda, \mu \neq 0$, क्योंकि $p =
\lambda\mu \neq 0$)। $\bigl((\lambda^n), (\mu^n)\bigr)$ की स्वतंत्रता: $n =
0, 1$ पर कोई संबंध $c + c' = 0$, $c\lambda + c'\mu = 0$ देता है, अतः
$c(\lambda - \mu) = 0$: $c = c' = 0$। विमा $2$ में दो \omterm{def:b1:vspaces:free}{स्वतंत्र} सदिश: अतः
\omterm{def:b1:vspaces:free}{आधार}। दुहरा मूल: $\bigl((n\lambda^n)\bigr)$ एक हल है, क्योंकि $s = 2\lambda$
के साथ $p = \lambda^2$:
\[
s(n+1)\lambda^{n+1} - p\,n\lambda^n
= \lambda^{n+2}\bigl(2(n+1) - n\bigr) = (n+2)\lambda^{n+2} ;
\]
$n = 0, 1$ पर स्वतंत्रता: $c = 0$, फिर $\lambda \neq 0$ के साथ $c'\lambda =
0$।

\textbf{14.} $D = X^2 - X - 6 = (X - 3)(X + 2)$। व्यापक हल $u_n = A\,3^n +
B(-2)^n$; आरंभिक शर्तें $A + B = 1$ और $3A - 2B = 8$ देती हैं, अतः $A = 2$,
$B = -1$:
\[
u_n = 2\cdot 3^n - (-2)^n .
\]
जाँच: $u_2 = 18 - 4 = 14 = u_1 + 6u_0$; $u_3 = 54 + 8 = 62 = u_2 + 6u_1 = 14
+ 48$।

\textbf{15.} $C\begin{pmatrix} u_n\\ u_{n+1}\end{pmatrix} = \begin{pmatrix}
u_{n+1}\\ -p\,u_n + s\,u_{n+1}\end{pmatrix} = \begin{pmatrix} u_{n+1}\\
u_{n+2}\end{pmatrix}$, और आगमन $C^n$ के साथ सूत्र दे देता है। इसके अतिरिक्त
$\operatorname{tr} C = 0 + s = s$ और $p(C) = 0\cdot s - 1\cdot(-p) = p$:
अर्थात् सहचर आव्यूह का केली--हैमिल्टन \omterm{def:b1:poly:def}{बहुपद} ठीक $D$ है।

\textbf{16.} $u_{n+3} = \alpha u_{n+2} + \beta u_{n+1} + \gamma u_n$ के किसी
भी हल $u$ के लिए अवस्था-सदिश $v_n = (u_n, u_{n+1}, u_{n+2})^{\mathsf T}$
$C_3 v_n = v_{n+1}$ संतुष्ट करते हैं (पहली दो पंक्तियाँ सरका देती हैं, और
अंतिम पंक्ति पुनरावृत्ति लगाती है)। अतः
\[
D_3(C_3)\,v_0 = v_3 - \alpha v_2 - \beta v_1 - \gamma v_0 ,
\]
जिसके तीनों घटक $u_{k+3} - \alpha u_{k+2} - \beta u_{k+1} - \gamma u_k = 0$
हैं ($k = 0, 1, 2$)। और चूँकि आरंभिक अवस्था $v_0 = (u_0, u_1, u_2)^{\mathsf
T}$ पूरे $K^3$ पर घूमती है (आरंभिक मान मुक्त हैं), अतः आव्यूह $D_3(C_3)$ हर
सदिश को मार देता है: $D_3(C_3) = 0$।

\textbf{17.} $\deg R_n \leq 2$ के साथ $X^n = Q\,D_3 + R_n$ लिखिए और हर मूल
पर मूल्यांकन कीजिए: $\lambda_i^n = R_n(\lambda_i)$। अतः $R_n$ $\leq 2$ घात
का ऐसा \omterm{def:b1:poly:def}{बहुपद} है जो तीन भिन्न गाँठों $\lambda_i$ पर तीनों मानों $\lambda_i^n$
का अंतर्वेशन करता है: \cref{thm:b1:poly:lagrange} की अद्वितीयता से $R_n =
\sum_i \lambda_i^n L_i$, जहाँ $(L_i)$ उन गाँठों का लाग्रांज \omterm{def:b1:vspaces:free}{आधार} है। $C_3$
प्रतिस्थापित करने पर (प्रश्न 5 और 16):
\[
C_3^{\,n} = R_n(C_3) = \sum_{i=1}^{3} \lambda_i^n\,L_i(C_3),
\]
जहाँ तीनों आव्यूह $L_i(C_3)$ $n$ से स्वतंत्र हैं: अतः $C_3^{\,n}$ की हर
प्रविष्टि $\lambda_1^n, \lambda_2^n, \lambda_3^n$ का कोई नियत संयोजन है।

\textbf{18.} $D_3 = X^3 - 2X^2 - X + 2 = (X-1)(X+1)(X-2)$। व्यापक हल $u_n =
A + B(-1)^n + C\,2^n$। आरंभिक शर्तें: $A + B + C = 0$, $A - B + 2C = 1$, $A
+ B + 4C = 1$। तीसरी में से पहली घटाने पर: $3C = 1$, $C = \frac13$; फिर $A +
B = -\frac13$ और $A - B = \frac13$: $A = 0$, $B = -\frac13$। इस प्रकार
\[
u_n = \frac{2^n - (-1)^n}{3}
\]
(याकोब्स्टाल अंक)। जाँच: $u_3 = \frac{8 + 1}{3} = 3 = 2u_2 + u_1 - 2u_0 = 2
+ 1 - 0$।

\textbf{19.} $\lambda$ पर \omterm{def:b1:poly:def}{बहुपद} $X^n$ का टेलर प्रसार:
\[
X^n = \sum_{k=0}^{n} \binom nk \lambda^{n-k}(X - \lambda)^k ,
\]
और $k \geq 3$ वाले सभी पद $(X - \lambda)^3$ से विभाज्य हैं: अतः शेषफल है
\[
R_n = \lambda^n + n\lambda^{n-1}(X - \lambda) + \binom
n2\lambda^{n-2}(X - \lambda)^2 .
\]
$N^3 = 0$ के साथ $M = \lambda I + N$ के लिए: $(M - \lambda I)^3 = N^3 = 0$,
अतः प्रश्न 5 देता है
\[
M^n = \lambda^n I + n\lambda^{n-1} N + \binom n2
\lambda^{n-2} N^2 ,
\]
जो ठीक $(\lambda I + N)^n$ का वही द्विपद प्रसार है, $N^2$ पर कटा हुआ ---
अर्थात् दोनों विधियाँ सहमत हैं।

\textbf{20.} हल-समष्टि की विमा $3$ है (प्रश्न 11 जैसा ही $(u_0, u_1, u_2)$
से प्राचलन), और हर $(\lambda_i^n)$ एक हल है। स्वतंत्रता: मान लीजिए $n = 0,
1, 2$ के लिए $c_1\lambda_1^n + c_2\lambda_2^n + c_3\lambda_3^n = 0$। $i$
नियत कीजिए और $L_i = \sum_{k \leq 2} p_k X^k$ को उन गाँठों का वह लाग्रांज
\omterm{def:b1:poly:def}{बहुपद} लीजिए जिसके लिए $L_i(\lambda_j) = \delta_{ij}$। तब
\[
0 = \sum_{k=0}^{2} p_k\Bigl(\sum_j c_j\lambda_j^k\Bigr)
= \sum_j c_j\,L_i(\lambda_j) = c_i .
\]
अतः सभी $c_i = 0$: विमा $3$ में तीन \omterm{def:b1:vspaces:free}{स्वतंत्र} हल, यानी एक \omterm{def:b1:vspaces:free}{आधार}; और व्यापक हल
$c_1\lambda_1^n + c_2\lambda_2^n + c_3\lambda_3^n$ है।

\textbf{21.} $F_k = F_{k+2} - F_{k+1}$ से योग दूरबीनी हो जाता है:
\[
\sum_{k=1}^{n} F_k = \sum_{k=1}^{n}\bigl(F_{k+2} - F_{k+1}\bigr)
= F_{n+2} - F_2 = F_{n+2} - 1 .
\]

\textbf{22.} $F^{m+n} = F^m F^n$ की $(1,2)$ प्रविष्टि लीजिए: बायाँ पक्ष
$F_{m+n}$ है; और दायाँ पक्ष ($F^m$ की $1$ पंक्ति) गुणा ($F^n$ का $2$ स्तंभ)
है, अर्थात् $F_{m+1}F_n + F_m F_{n-1}$। $m = n$ के साथ:
\[
F_{2n} = F_{n+1}F_n + F_nF_{n-1} = F_n\,(F_{n+1} + F_{n-1}).
\]

\textbf{23.} बिने से $F_n - \dfrac{\varphi^n}{\sqrt5} =
-\dfrac{\psi^n}{\sqrt5}$, और $\abs\psi = \frac{\sqrt5 - 1}2 < 1$, अतः
\[
\Bigl|F_n - \frac{\varphi^n}{\sqrt5}\Bigr|
\leq \frac{1}{\sqrt5} < \frac12
\qquad (n \geq 0):
\]
$F_n$ $\varphi^n/\sqrt5$ के सबसे निकट का पूर्णांक है।

\textbf{24.} $t_n = F_{n+1} + F_{n-1}$ सरकाए हुए फ़िबोनाच्ची अनुक्रमों का
संयोजन है, अतः वह वही पुनरावृत्ति संतुष्ट करता है: $t_{n+2} = t_{n+1} +
t_n$; और $t_1 = F_2 + F_0 = 1$, $t_2 = F_3 + F_1 = 3$: ये लूका अंक $L_n$
हैं। अनुक्रम $\varphi^n + \psi^n$ ऐसा हल है जिसके पहले दो मान वही हैं
($\varphi + \psi = 1$, $\varphi^2 + \psi^2 = ( \varphi + \psi)^2 -
2\varphi\psi = 3$), अतः $L_n = \varphi^n + \psi^n$। अंत में
\[
F_n L_n = \frac{(\varphi^n - \psi^n)(\varphi^n +
\psi^n)}{\sqrt5} = \frac{\varphi^{2n} - \psi^{2n}}{\sqrt5} =
F_{2n},
\]
जिससे प्रश्न 22 फिर से मिल जाता है।

\textbf{25.} (क) पाँचों आव्यूह $I, A, A^2, A^3, A^4$ $4$-विमीय
$\mathcal{M}_2(K)$ में रहते हैं, अतः $\leq 4$ घात का \emph{कोई न कोई} अशून्य
\omterm{def:b1:poly:def}{बहुपद} $A$ को मार देता है; और भाग II ने इसे स्पष्ट द्विघातीय $A^2 = sA - pI$
तक तीक्ष्ण कर दिया, जो सभी घातों को समतल $\operatorname{Vect}(I, A)$ में बंद
कर देता है। (ख) यूक्लिडीय भाग $X^n$ का उस द्विघातीय के सापेक्ष लघुकरण कर
देता है, और शेषफल के दोनों गुणांक दो-पदीय पुनरावृत्ति $a_{n+2} = s\,a_{n+1}
- p\,a_n$ का पालन करते हैं: इस प्रकार घातांकन पुनरावृत्ति बन गया। (ग) प्रश्न
6 \emph{केली--हैमिल्टन प्रमेय} की $n = 2$ वाली स्थिति है, जो हर विमा में
मान्य है और स्नातक वर्ष~2 के खंड में सिद्ध की गई है। (घ) सहचर आव्यूह घेरा
पूरा कर देता है: हर \omterm{pb:b1:matrices:1}{रैखिक पुनरावृत्ति} किसी आव्यूह की घात \emph{है}, और वही
\omterm{def:b1:poly:def}{बहुपद} $D$ अनुरेख-और-सारणिक के आँकड़ों के रूप में आ जाता है, अतः शेषफल-कलन एक
ही झटके में पुनरावृत्तियाँ हल करता है और घातें संगणित करता है।
\end{solution}
